\scanpage{159}
Notons que les topos du type (258) se ramènent au type (257), en remplaçant
\[
X,\ \deleted{\unsure{\Sigma}},\Sigma',\Sigma'',\Sigma\quad\text{par}\quad
\deleted{\unsure{X}}X\setminus\Sigma'_0,\ \Sigma'\setminus\Sigma'_0,\ \Sigma''\setminus\Sigma'_0,\ \Sigma\setminus\Sigma\cap\Sigma'_0.
\]
Pour faire cette \unsure{déduction|réduction}, il faut dans notre formalisme savoir expliciter \margin{$\Vert$} le passage au \deleted{\unsure{X}} \uline{complémentaire} d'une partie admissible $\Sigma'_0$. Il faudra y revenir plus bas. D'autre part, le cas $V_{\unsure{\Sigma'},\Sigma''}^{\Sigma}$ se ramène au cas $V_{\Sigma',X}^{\Sigma}$ (où $\Sigma''=X$), pourvu qu'on sache expliciter le passage de $X$ à un $\Sigma''$ -- bien entendu, on perd alors l'hypothèse des croisements normaux -- qui, à dire vrai, \unsure{elle} ne joue plus de rôle depuis belle lurette!\ldots) -- c'est \unsure{fait}, \unsure{sûrement}, par le système des $D_{d_0\ldots d_r}^{\Sigma''}$. \deleted{\unsure{Enfin}}

En fait, au niveau des topos, il
\margin{158}
